\documentclass[11pt]{article}
% Shared preamble for Vietnamese companion manuscripts
\usepackage{fontspec}
\setmainfont{DejaVu Serif}
\usepackage[a4paper,margin=2.1cm]{geometry}
\usepackage{microtype}
\usepackage{amsmath,amssymb}
\usepackage{booktabs,tabularx,array,longtable}
\usepackage{graphicx}
\usepackage{enumitem}
\usepackage[hidelinks]{hyperref}
\usepackage{xurl}
\usepackage{titlesec}
\usepackage{fancyhdr}
\setlength{\parindent}{0pt}
\setlength{\parskip}{0.55em}
\setlist[itemize]{leftmargin=1.5em,itemsep=0.25em}
\setlist[enumerate]{leftmargin=1.7em,itemsep=0.25em}
\renewcommand{\arraystretch}{1.12}
\titleformat{\section}{\large\bfseries}{\thesection}{0.6em}{}
\titleformat{\subsection}{\normalsize\bfseries}{\thesubsection}{0.5em}{}
\pagestyle{fancy}\fancyhf{}\fancyfoot[C]{\small Bản tiếng Việt đồng hành -- trang \thepage}\renewcommand{\headrulewidth}{0pt}

\emergencystretch=3em
\sloppy

\title{\textbf{GovRed-Health: Đánh giá sai lệch quyền và quản trị hai chiều thời gian giữa thời điểm cung cấp dữ liệu và ghi bền vững}}
\author{Nguyễn Ngọc Thiện, Trịnh Minh Quang \\ Trường THPT Hai Bà Trưng, Huế, Việt Nam \\
Tác giả liên hệ: \texttt{kakangocthien109@gmail.com}}
\date{2026}
\begin{document}
\maketitle

\begin{abstract}
Các hệ thống AI y tế có trạng thái vận hành trên các nền tảng dữ liệu lớn sức khỏe dài hạn thường xuyên truy vấn bối cảnh bệnh án lịch sử và sau đó đề xuất các sửa đổi trạng thái bền vững. Tuy nhiên, trong quy trình lâm sàng phân tán, trạng thái người bệnh, chân trời hai chiều thời gian (bitemporal: thời gian hiệu lực và thời gian tri thức), sự đồng thuận động, chính sách quản trị, quyền hạn vai trò, ràng buộc mục đích, trạng thái bộ nhớ đệm hoặc hiệu lực ảnh chụp ngữ cảnh có thể bị sai lệch trong khoảng thời gian từ lúc cung cấp dữ liệu đến khi ghi bền vững ($\Delta t = t_{\text{commit}} - t_{\text{disc}}$). Dựa trên nền tảng hệ thống dữ liệu hai chiều thời gian kinh điển của Snodgrass, GovRed-Health thiết lập một bộ đánh giá thực chứng về hiện tượng sai lệch quyền (authorization drift) trong môi trường dữ liệu lớn y tế có trạng thái. Mỗi ca thử nghiệm là một lịch logic ghép cặp: hệ thống cung cấp cho mô hình một bối cảnh sức khỏe có quản trị giới hạn theo tác vụ, chủ động kích hoạt các đột biến trạng thái bitemporal hoặc điều kiện quản trị, và kiểm tra tính bền vững xuyên qua ứng dụng, lớp phân quyền, bộ nhớ đệm Redis, PostgreSQL và ranh giới kiểm toán mật mã. Thiết kế đánh giá cô lập các cơ chế quản trị dài hạn có khả năng mở rộng đồng thời cam đoan không rò rỉ thông tin y tế được bảo vệ (zero-PHI leakage) thông qua cơ chế kiểm tra token sentinel ngẫu nhiên độ hỗn loạn cao. Trong đợt chạy niêm phong \texttt{2026-08-17-rivf-final-003} ($N=270$ lịch tổng hợp mỗi nhánh, 210 lịch chính có thể thực thi), tỷ lệ điểm cuối tổng hợp không an toàn đã ấn định trước lần lượt là 57,1\% (120/210; 95\% KTC: 0,504--0,636) ở Unbound, 42,9\% (90/210; 95\% KTC: 0,364--0,496) ở State-Version-Only, 42,9\% (90/210; 95\% KTC: 0,364--0,496) ở Snapshot-Bound-State-Only và 14,3\% (30/210; 95\% KTC: 0,102--0,197) ở GLHS-Strict. Kiểm định McNemar chính xác khẳng định ưu thế vượt trội của Strict so với Unbound ($b=90, c=0, p=1{,}62\times 10^{-27}$). Toàn bộ 30 ca thất bại còn lại của Strict đều thuộc nhóm ghi trên trạng thái cũ trong điều kiện đồng thời có trật tự tuyến tính hóa chưa phân định; các nhóm sai lệch quyền tuần tự không tạo ra bất kỳ thao tác ghi sai nào ở nhánh Strict. Tỷ lệ cung cấp dữ liệu bị cấm và rò rỉ PHI đạt mức tuyệt đối 0/270 (0,00\%) ở tất cả các nhánh.
\end{abstract}

\textbf{Từ khóa:} dữ liệu lớn y tế, hệ thống dữ liệu bitemporal, AI y tế, sai lệch quyền, TOCTOU, quản trị dài hạn, zero-PHI leakage, tác tử có trạng thái, benchmark.

\section{Giới thiệu}
Các nền tảng dữ liệu lớn y tế dài hạn tích hợp các luồng sự kiện lâm sàng append-only từ nhiều nguồn phân tán: bệnh án điện tử (EHR), thiết bị giám sát từ xa và kho xét nghiệm. Trong các kiến trúc AI y tế có trạng thái, các tác tử mô hình ngôn ngữ lớn (LLM) liên tục đọc dữ liệu lịch sử bệnh án để đưa ra khuyến nghị điều trị, soạn thảo tóm tắt lâm sàng hoặc tạo ra các cập nhật trạng thái bền vững. Mô hình tương tác phân tách này làm nảy sinh một lỗ hổng toàn vẹn xuyên thao tác cốt lõi: một lần cung cấp dữ liệu có thể hoàn toàn hợp lệ khi chuyển giao cho mô hình tại thời điểm suy luận $t_{\text{disc}}$, nhưng lại trở thành một căn cứ trái phép hoặc sai lệch cho thao tác ghi bền vững tại thời điểm xác nhận $t_{\text{commit}}$ sau khi trạng thái người bệnh, chứng cứ lâm sàng, sự đồng thuận của bệnh nhân, chính sách y tế, vai trò người dùng hoặc ảnh chụp ngữ cảnh đã thay đổi.

Hiện tượng sai lệch quyền này là một dạng lỗi Time-of-Check to Time-of-Use (TOCTOU) đặc thù, khác biệt với các cuộc tấn công đơn lẻ như prompt injection hay lỗi truy xuất thông tin đơn thuần. Các khung đánh giá bảo mật hiện nay chủ yếu tập trung vào tấn công bẻ khóa đối kháng (AgentDojo~\cite{agentdojo}, ASB~\cite{asb}), prompt injection trong y khoa (MPIB~\cite{mpib}) hoặc cô lập luồng đọc dữ liệu nhạy cảm (SecureClaw~\cite{secureclaw}). Song song đó, các nghiên cứu về giao dịch tác tử đã đề xuất kiểm tra quyền ở thời điểm ghi (CommitGuard~\cite{commitguard}), khả năng tuần tự hóa trạng thái chính sách (Provenact~\cite{provenact}), kiểm soát truy cập bộ nhớ đa chủ thể (GateMem~\cite{gatemem}) và đại số bộ nhớ hai chiều thời gian (TOKI~\cite{toki}). Bảng~\ref{tab:benchmark_comparison_vi} so sánh năng lực bao phủ các chiều đe dọa giữa các bộ đánh giá hiện hành.

\begin{table}[H]
\caption{So sánh năng lực bao phủ các chiều đe dọa và cơ chế hệ thống giữa các bộ đánh giá bảo mật tác tử y tế.}
\label{tab:benchmark_comparison_vi}
\centering\small
\setlength{\tabcolsep}{4pt}
\begin{tabularx}{\textwidth}{l c c c c c c}
\toprule
\textbf{Bộ đánh giá} & \textbf{Tấn công Prompt} & \textbf{Quyền Công cụ} & \textbf{Lệch quyền TOCTOU} & \textbf{Thu hồi Đồng thuận} & \textbf{Bitemporal} & \textbf{Dấu vết Merkle}\\
\midrule
AgentDojo~\cite{agentdojo} & \checkmark & \checkmark & $\times$ & $\times$ & $\times$ & $\times$\\
MPIB~\cite{mpib} & \checkmark & $\times$ & $\times$ & $\times$ & $\times$ & $\times$\\
CommitGuard~\cite{commitguard} & $\times$ & \checkmark & \checkmark & $\times$ & $\times$ & $\times$\\
Provenact~\cite{provenact} & $\times$ & \checkmark & \checkmark & $\times$ & $\times$ & $\times$\\
GateMem~\cite{gatemem} & $\times$ & \checkmark & $\times$ & $\times$ & $\times$ & $\times$\\
TOKI~\cite{toki} & $\times$ & \checkmark & $\times$ & $\times$ & \checkmark & $\times$\\
\textbf{GovRed-Health (Chúng tôi)} & \checkmark & \checkmark & \textbf{\checkmark} & \textbf{\checkmark} & \textbf{\checkmark} & \textbf{\checkmark}\\
\bottomrule
\end{tabularx}
\end{table}

GovRed-Health giải quyết khoảng trống quản trị hệ thống cụ thể trong dữ liệu lớn y tế: \textbf{điều gì xảy ra sau khi mô hình đã thực sự sử dụng một phần dữ liệu sức khỏe được cung cấp đúng quyền, và trạng thái dài hạn, dòng thời gian bitemporal hoặc thẩm quyền quản trị thay đổi trước khi thao tác ghi được thực hiện?} Bộ đánh giá tích hợp lý thuyết hệ thống dữ liệu hai chiều thời gian kinh điển của Richard T. Snodgrass~\cite{snodgrass1995,kumar1998}, phân biệt rạch ròi giữa thời gian hiệu lực (khi một khẳng định lâm sàng thực sự đúng trong thực tế) và thời gian giao dịch/tri thức (khi khẳng định được ghi nhận vào cơ sở dữ liệu). Bằng cách nhúng tọa độ bitemporal và niêm phong mật mã Merkle vào đối tượng cung cấp dữ liệu, GovRed-Health đánh giá khả năng quản trị dữ liệu dài hạn có khả năng mở rộng với cam đoan không rò rỉ thông tin y tế nhạy cảm (zero-PHI leakage).

Đóng góp của nghiên cứu gồm: (1) trừu tượng hóa lịch logic có đóng băng bối cảnh bitemporal mà AI đã nhìn thấy vào đơn vị phân tích thống kê; (2) bộ đột biến trạng thái và quản trị ghép cặp trải rộng trên các chiều chủ thể, đồng thuận, chính sách, vai trò/mục đích, trạng thái bitemporal, bộ nhớ đệm và kiểm toán; (3) thực thi ranh giới thực xuyên qua PostgreSQL, Redis và nhật ký kiểm toán mật mã; và (4) giao thức phân tích tách biệt rạch ròi giữa từ chối phòng vệ an toàn, lỗi hạ tầng và điều kiện đồng thời chưa phân định thứ tự dưới cam đoan zero-PHI leakage.

\section{Nền tảng bitemporal và mô hình đe dọa}
\subsection{Ranh giới tính mới}
GovRed-Health không tự nhận phát minh ra prompt injection, lý thuyết kiểm soát truy cập, đại số cơ sở dữ liệu bitemporal, kiểm soát đồng thời lạc quan hay kiểm tra quyền tại thời điểm ghi. Bốn cấu hình hệ thống được đánh giá là các phép loại bỏ thành phần (ablation) nội bộ nhằm cô lập đóng góp của từng cơ chế kỹ thuật, không phải là bản cài đặt chuẩn của các hệ thống đã công bố trước đây.

\subsection{Nền tảng hai chiều thời gian Snodgrass}
Kho dữ liệu lớn lâm sàng không thể chỉ dựa vào một mốc thời gian vật lý duy nhất. Theo nền tảng cơ sở dữ liệu thời gian của Snodgrass~\cite{snodgrass1995}, mỗi khẳng định lâm sàng $\mathcal{E}$ được đặc trưng bởi hai chiều thời gian trực giao:
\begin{enumerate}
\item \textbf{Thời gian hiệu lực (Valid Time)} ($\tau_v = [v_{\text{start}}, v_{\text{end}})$): Khoảng thời gian liên tục mà một trạng thái sinh lý, bệnh lý hoặc phác đồ dùng thuốc thực sự tồn tại trong cơ thể người bệnh.
\item \textbf{Thời gian giao dịch / tri thức (Transaction / Knowledge Time)} ($\tau_k = [t_{\text{start}}, t_{\text{end}})$): Khoảng thời gian rời rạc mà khẳng định lâm sàng được hệ thống cơ sở dữ liệu ghi nhận chính thức và chưa bị sửa đổi hay thay thế.
\end{enumerate}
Các hiệu chỉnh bệnh án hồi cứu, nhập liệu chẩn đoán muộn hoặc kết quả xét nghiệm gửi trễ sẽ làm biến đổi $\tau_k$ hoặc $\tau_v$. Khi tác tử AI thực hiện suy luận đa bước trong khoảng thời gian $\Delta t = t_{\text{commit}} - t_{\text{disc}}$, các sửa đổi hồi cứu này tạo ra sự sai lệch bitemporal, khiến cơ sở tiền đề của AI bị vô hiệu ngay cả khi phiên bản trạng thái tức thời có vẻ mới.

\subsection{Đối tượng cung cấp dữ liệu có quản trị}
Một lần cung cấp dữ liệu có quản trị bitemporal được biểu diễn bởi bộ:
\begin{equation}
H = (u, a, p, \tau_v, \tau_k, v_s, v_g, d, e, E),
\end{equation}
trong đó $u$ là chủ thể người bệnh, $a$ là tác nhân được xác thực, $p$ là mục đích lâm sàng được phê duyệt, $\tau_v$ là chân trời thời gian hiệu lực, $\tau_k$ là chân trời thời gian tri thức, $v_s$ là tọa độ phiên bản trạng thái phân vùng thực thể, $v_g = (v_c, v_p)$ là tọa độ đồng phiên bản gồm đồng thuận ($v_c$) và chính sách ($v_p$), $d = \mathcal{H}(E)$ là mã băm mật mã Merkle chuẩn tắc, $e$ là thời hạn hiệu lực của hợp đồng thuê (lease), và $E$ là tập bằng chứng lâm sàng được cấp phép theo tác vụ.

Một đề xuất ghi bền vững $P(H, \Delta)$ yêu cầu một biến đổi trạng thái $\Delta$. GovRed-Health làm đột biến chính xác một phụ thuộc trong $H$ trong khoảng $\Delta t$ trong khi vẫn giữ nguyên hình dạng cấu trúc của tác vụ.

\subsection{Mô hình đe dọa và cam đoan Zero-PHI Leakage}
Chúng tôi mô hình hóa các yêu cầu API đã được xác thực, bộ nhớ đệm khách bị lỗi thời, các luồng ghi đồng thời, việc sửa đổi bệnh án bitemporal hồi cứu, sự thu hồi đồng thuận của bệnh nhân theo chuẩn HIPAA/GDPR~\cite{fhirconsent,fhirprov} và các đoạn văn bản chứng cứ chưa qua kiểm duyệt. Mô hình đe dọa loại trừ việc kẻ tấn công chiếm quyền root cơ sở dữ liệu hoặc đánh cắp khóa bí mật hạ tầng.

Để đánh giá tính riêng tư trên dữ liệu lớn y tế mà không làm lộ thông tin nhạy cảm, toàn bộ thử nghiệm được thực thi trên đoàn hệ bệnh nhân tổng hợp có nhúng các token sentinel canary ngẫu nhiên với độ hỗn loạn cao $\sigma_k \in \{0, 1\}^{256}$ tại các trường thông tin không được phép tiết lộ. Bất biến zero-PHI leakage yêu cầu:
\begin{equation}
\forall H, \quad \operatorname{PHI\_Leak}(H) = 0 \quad \land \quad \operatorname{Sentinel\_Leak}(H, \text{Audit}) = 0.
\end{equation}

\section{Thiết kế đánh giá}
\subsection{Các lịch logic ghép cặp}
Bảng~\ref{tab:schedules_vi} xác định các họ lịch kiểm thử sai lệch quyền cốt lõi. Đơn vị phân tích thống kê là một lịch logic hoàn chỉnh ($H \to \text{Đột biến} \to P \to \text{Kiểm toán}$), không phải là một lệnh gọi mạng riêng lẻ.

\begin{table}[H]
\caption{Các họ lịch kiểm thử sai lệch quyền và quản trị bitemporal.}
\label{tab:schedules_vi}
\centering\small
\setlength{\tabcolsep}{4pt}
\begin{tabularx}{\textwidth}{p{0.36\textwidth}X}
\toprule
\textbf{Họ lịch logic} & \textbf{Bất biến hệ thống chính} \\
\midrule
Cung cấp dữ liệu $\to$ thu hồi đồng thuận $\to$ ghi & Đồng thuận bị thu hồi lập tức vô hiệu hóa các đề xuất ghi phát sinh từ lần đọc trước đó. \\
Cung cấp dữ liệu $\to$ cập nhật chính sách $\to$ ghi & Sai lệch chính sách được tái thẩm định nghiêm ngặt tại thời điểm ghi bền vững. \\
Cung cấp dữ liệu $\to$ đổi tác nhân/vai trò/mục đích & Tọa độ ủy quyền không thể tự động gắn lại sang vai trò hoặc mục đích khác. \\
Phát lại chéo chủ thể & Ảnh chụp ngữ cảnh và đề xuất không thể hoán đổi giữa các bệnh nhân khác nhau. \\
Tiến triển trạng thái bitemporal $\to$ ghi trên bản cũ & Cập nhật hồi cứu $\tau_v/\tau_k$ vô hiệu hóa các thao tác ghi dựa trên phiên bản cơ sở đã lỗi thời. \\
Thu hồi quyền $\to$ tái sử dụng bộ nhớ đệm & Bộ nhớ đệm phụ Redis không được phép lưu trữ hoặc cung cấp dữ liệu đã bị thu hồi quyền. \\
Hết hạn / Sai mã băm Merkle & Hợp đồng thuê bị hết hạn hoặc dữ liệu chứng cứ bị sửa đổi sẽ bị từ chối ghi. \\
Tiến trình quản trị $\parallel$ Đề xuất ghi đồng thời & Kết quả phải tuyến tính hóa chính xác theo trạng thái được xác nhận trước; bất định nếu tranh chấp. \\
Tái dựng kiểm toán bitemporal & Toàn bộ dòng dõi ($H \to \Delta \to \text{Quyết định} \to \text{Ghi}$) phải tái dựng được với 0\% lộ PHI. \\
\bottomrule
\end{tabularx}
\end{table}

\subsection{Các nhánh thực nghiệm}
Bốn nhánh thử nghiệm được cấu hình trên cùng một đường dẫn ứng dụng:
\begin{itemize}
\item \textbf{UNBOUND}: Chấp thuận ghi đề xuất tự do mà không ràng buộc với ảnh chụp hay chân trời thời gian.
\item \textbf{STATE\_VERSION\_ONLY}: Chỉ kiểm tra phiên bản trạng thái thực thể $v_s$, bỏ qua mã băm ảnh chụp và hợp đồng thuê bitemporal.
\item \textbf{SNAPSHOT\_BOUND\_STATE\_ONLY}: Ràng buộc đề xuất với mã băm ảnh chụp $d$ và trạng thái $v_s$, nhưng không kiểm tra lại chính sách đồng thuận động $v_g$ tại thời điểm ghi.
\item \textbf{GLHS\_STRICT}: Thực thi toàn diện tính nhất quán khoảng bitemporal Snodgrass ($\tau_v \cap \tau_k$), ràng buộc mã băm Merkle chính xác ($d = \mathcal{H}(E)$) và tái thẩm định thẩm quyền quản trị động ($v_g$).
\end{itemize}

\subsection{Thực thi trên ranh giới thực và khả năng mở rộng}
Bộ đánh giá được thực thi xuyên qua ranh giới ứng dụng thực tế với cơ sở dữ liệu PostgreSQL, lớp đệm Redis và bảng kiểm toán ghi nối tiếp bất biến. Để đảm bảo khả năng mở rộng cho dữ liệu lớn y tế, cơ chế kiểm soát đồng thời lạc quan (OCC) được phân vùng theo đồ thị phụ thuộc thực thể (Entity DAG), loại bỏ triệt để hiện tượng suy thoái hiệu năng do nghẽn khóa toàn bảng (hiện tượng Thomasian thrashing). Các bộ quan sát đo lường tuân thủ nguyên tắc read-only tuyệt đối.

\subsection{Bộ thẩm định tất định}
Các bộ thẩm định tự động (deterministic oracles) chấm điểm các thuộc tính ranh giới quan sát được mà không phụ thuộc vào việc LLM tự đánh giá. Bộ thẩm định kiểm tra việc lộ dữ liệu chéo chủ thể, rò rỉ token sentinel, chấp thuận ghi sai phiên bản, vi phạm khoảng thời gian bitemporal, sai lệch mã băm và khả năng tái dựng kiểm toán Merkle đầy đủ.

\section{Kế hoạch phân tích}
Đối với mỗi nhánh tương thích, báo cáo công bố tử số, mẫu số và khoảng tin cậy Wilson 95\% cho chỉ số tổng hợp không an toàn và tỷ lệ cung cấp dữ liệu bị cấm. Các thao tác ghi sai đã xác nhận, từ chối phòng vệ an toàn và trật tự đồng thời bất định được phân loại thành các lớp kết quả riêng biệt. So sánh ghép cặp giữa các nhánh sử dụng kiểm định McNemar hai phía chính xác với hiệu chỉnh Holm. Các chỉ số thứ cấp đánh giá sự vô hiệu hóa bộ nhớ đệm, khả năng tái dựng kiểm toán, độ trễ p50/p95 và mức độ bảo đảm zero-PHI leakage.

\section{Kết quả}
\subsection{Kết quả ranh giới chính}
Đợt chạy niêm phong \texttt{2026-08-17-rivf-final-003} (mã nguồn \texttt{5b2c0dbf17e2}) đã thực thi 270 lịch logic tổng hợp cho mỗi nhánh. Tập phân tích chính gồm 210 lịch có thể thực thi ở mỗi nhánh; 180 hàng bị loại trừ theo giao thức vẫn giữ trạng thái \textsc{Not Run} và không tính vào mẫu số.

\begin{table}[H]
\caption{Kết quả ranh giới trong đợt chạy niêm phong. Cột đầu tiên là chỉ số tổng hợp không an toàn đã ấn định trước. Toàn bộ 30 ca dư của Strict là các cuộc đua đồng thời có trật tự chưa phân định. Tỷ lệ cung cấp dữ liệu bị cấm đạt mức 0 tuyệt đối (zero-PHI leakage).}
\label{tab:results_vi}
\centering\small
\setlength{\tabcolsep}{6pt}
\begin{tabularx}{\textwidth}{l r r}
\toprule
\textbf{Nhánh thử nghiệm} & \textbf{Chỉ số tổng hợp không an toàn (95\% KTC)} & \textbf{Cung cấp dữ liệu bị cấm} \\
\midrule
UNBOUND & 57,1\% (120/210; 0,504--0,636) & 0/270 (0,0\%) \\
STATE\_VERSION\_ONLY & 42,9\% (90/210; 0,364--0,496) & 0/270 (0,0\%) \\
SNAPSHOT\_BOUND\_STATE\_ONLY & 42,9\% (90/210; 0,364--0,496) & 0/270 (0,0\%) \\
GLHS\_STRICT & 14,3\% (30/210; 0,102--0,197) & 0/270 (0,0\%) \\
\bottomrule
\end{tabularx}
\end{table}

Như trình bày tại Bảng~\ref{tab:results_vi}, GLHS-Strict đã giảm tỷ lệ không an toàn từ 57,1\% (Unbound) và 42,9\% (State-Version-Only) xuống còn 14,3\%. Kiểm định McNemar ghép cặp chính xác giữa GLHS-Strict và Unbound ghi nhận $b=90$ cặp bất đồng nghiêng về Strict, $c=0$ nghiêng về Unbound, với mức ý nghĩa thống kê vượt trội $p=1{,}62\times 10^{-27}$.

Đặc biệt, toàn bộ 30 ca thất bại còn lại của GLHS-Strict đều thuộc về các lịch đua ghi đồng thời trên trạng thái cũ, nơi trật tự xác nhận giao dịch trong cơ sở dữ liệu chưa thể phân định chắc chắn; chúng không cấu thành thao tác ghi sai đã xác nhận. Trong tất cả các họ sai lệch quyền tuần tự (thu hồi đồng thuận, đổi chính sách, đổi vai trò, sửa đổi bệnh án bitemporal hồi cứu, giả mạo ảnh chụp), GLHS-Strict đạt \textbf{0 thao tác ghi sai}.

\subsection{Quản trị mở rộng, kiểm toán bitemporal và Zero-PHI Leakage}
Các bộ quan sát vô hiệu hóa bộ nhớ đệm xác nhận tỷ lệ lỗi thu hồi bộ nhớ đệm là 0/270 ở mọi nhánh. Tỷ lệ cung cấp dữ liệu bị cấm và phát tán token sentinel canary đạt mức \textbf{0/270 (0,00\%)} ở tất cả các nhánh, chứng minh khả năng bảo toàn zero-PHI leakage trên thực nghiệm. Chuỗi kiểm toán bitemporal tái dựng thành công 100\% các chuyển đổi trạng thái đã xác nhận với đầy đủ tính toàn vẹn của thời gian hiệu lực ($\tau_v$), thời gian tri thức ($\tau_k$) và mã băm Merkle ($d$).

\section{Thảo luận}
GovRed-Health chứng minh rằng một phản hồi HTTP 200 thành công có thể che giấu những sai sót quản trị nghiêm trọng trong hệ thống dữ liệu lớn y tế: tác tử AI có thể thực hiện thành công một thao tác ghi bệnh án dựa trên bối cảnh đã bị vô hiệu hóa hồi cứu, sửa đổi hoặc thu hồi quyền.

Bằng cách tích hợp nền tảng dữ liệu hai chiều thời gian Snodgrass, nghiên cứu chỉ ra rằng việc chỉ kiểm tra phiên bản trạng thái $v_s$ là không đủ; hệ thống cần thẩm định tính nhất quán của giao điểm thời gian hiệu lực và thời gian tri thức ($\tau_v \cap \tau_k$) tại thời điểm ghi bền vững. Hơn nữa, việc phân tách giữa độ mới của trạng thái, tính toàn vẹn bitemporal và kiểm tra quyền động cho phép các nền tảng dữ liệu lớn y tế mở rộng quy mô thông qua OCC phân vùng thực thể mà không gặp rủi ro rò rỉ dữ liệu PHI hay sai lệch quyền thầm lặng.

\section{Giới hạn và đạo đức}
Bộ đánh giá hiện sử dụng dữ liệu bệnh nhân tổng hợp và một kiến trúc ứng dụng duy nhất; giá trị khái quát hóa sang các hệ thống EHR thương mại không đồng nhất cần được thẩm định đa trung tâm. Các oracle canary đo lường các vi phạm ranh giới quan sát được chứ không bao quát mọi tổn hại riêng tư ngữ nghĩa. Các nhánh hạ cấp chỉ phục vụ mục đích nghiên cứu. Nghiên cứu tập trung vào quản trị hệ thống phần mềm, không đánh giá hiệu quả chẩn đoán lâm sàng hay sự tuân thủ pháp lý.

\section{Khả năng tái lập}
Mọi đợt thực thi được ràng buộc với tập lịch logic đã đóng băng, nhật ký thực thi, cấu hình PostgreSQL/Redis và bảng mã băm SHA-256 dưới mã định danh bất biến \texttt{2026-08-17-rivf-final-003} (git commit \texttt{5b2c0dbf17e2}). Nhật ký định dạng máy đọc và mã nguồn phân tích được lưu trữ đầy đủ trong kho lưu trữ tạo phẩm.

\section{Kết luận}
GovRed-Health thiết lập một chuẩn đánh giá thực chứng về hiện tượng sai lệch quyền xuyên suốt khoảng thời gian từ cung cấp dữ liệu đến ghi bền vững trong hệ thống dữ liệu lớn y tế. Dựa trên nền tảng bitemporal của Snodgrass và cơ chế quản trị dài hạn có khả năng mở rộng, GLHS-Strict đã giảm chỉ số không an toàn từ 57,1\% xuống 14,3\% ($p=1{,}62\times 10^{-27}$), không ghi nhận bất kỳ thao tác ghi sai tuần tự nào và bảo toàn zero-PHI leakage (0/270) trên toàn bộ các nhánh. Kết quả này cung cấp nền tảng vững chắc cho việc xây dựng các kiến trúc AI y tế có trạng thái an toàn và có thể kiểm chứng.

\begin{thebibliography}{00}
\bibitem{snodgrass1995} R. T. Snodgrass, \emph{The TSQL2 Temporal Query Language}, Boston, MA: Kluwer Academic Publishers, 1995.
\bibitem{kumar1998} A. Kumar, V. J. Tsotras, and C. Faloutsos, ``Designing access methods for bitemporal databases,'' \emph{IEEE Trans. Knowl. Data Eng.}, vol. 10, no. 1, pp. 1--20, 1998.
\bibitem{agentdojo} E. Debenedetti, J. Zhang, M. Balunovi\'c, L. Beurer-Kellner, M. Fischer, and F. Tram\`er, ``AgentDojo: A dynamic environment to evaluate prompt injection attacks and defenses for LLM agents,'' in \emph{Advances in Neural Information Processing Systems}, Datasets and Benchmarks Track, 2024.
\bibitem{asb} H. Zhang, J. Huang, K. Mei, Y. Yao, Z. Wang, C. Zhan, H. Wang, and Y. Zhang, ``Agent Security Bench (ASB): Formalizing and benchmarking attacks and defenses in LLM-based agents,'' in \emph{International Conference on Learning Representations}, 2025.
\bibitem{mpib} J. Lee, H. Jang, and K. S. Choi, ``MPIB: A benchmark for medical prompt injection attacks and clinical safety in LLMs,'' arXiv:2602.06268, 2026.
\bibitem{secureclaw} Y. Ma and S. Schmid, ``SecureClaw: Clawing back control of LLM agents,'' arXiv:2606.09549, 2026.
\bibitem{commitguard} I. Santos-Grueiro, ``Temporary Authority, Permanent Effects: Commit-time authorization for LLM agents,'' arXiv:2607.10487, 2026.
\bibitem{provenact} Y. Peng and X. Wu, ``Stateful governance for concurrent agentic systems,'' arXiv:2608.02764, 2026.
\bibitem{gatemem} Z. Ren, Y. Yang, Y. Chen, et al., ``GateMem: Benchmarking memory governance in multi-principal shared-memory agents,'' arXiv:2606.18829, 2026.
\bibitem{toki} Z. Wang, ``TOKI: A bitemporal operator algebra for contradiction resolution in LLM-agent persistent memory,'' arXiv:2606.06240, 2026.
\bibitem{fhirconsent} HL7 International, ``FHIR R4: Consent,'' 2019.
\bibitem{fhirprov} HL7 International, ``FHIR R4: Provenance,'' 2019.
\end{thebibliography}

\end{document}
